\section{Basics}

\subsection{Grundsätzlich}

C++ ist sehr ähnlich zu C. Es gibt aber einige wichtige Änderungen zu C:

\begin{itemize}[itemsep=1pt, parsep=0pt]
    \item zu benutzender Compiler heisst jetzt \texttt{clang++}
    \item Nativen Bool Typ : \textbf{bool}
    \item richtiges Konstanten Schlüsselwort : \textbf{constexpr}
    \item Standartbibliotheken haben kein .h mehr beim Aufrufen
    \item C Standartbibliotheken können weiter verwendet werden, haben aber ein C im Header(stdio.h $\rightarrow$ cstdio)
    \item Char-Literale (z. B. 'A') haben in C++ den Datentyp char (in C war es int)(Brauchen ein const)
    \item Es gibt benannte Namensräume. Wichtigster Namensraum für die STL: \texttt{std}
\end{itemize}

\subsection{Hello world C++}

\lstinputlisting{code/01_basics/hello_World.cpp}

\subsection{Range-based for Loop}\label{forEach}

Range based for Loops laufen automatisch ein Array ab, ohne das zuerst die Grösse ermittelt werden muss. 
Die C for loop syntax wurde etwas erweitert.
Syntax:

\lstinputlisting{code/01_basics/forEach.cpp}

Siehe: Als Array Name wird ein \texttt{Pointer} auf das Array verlangt, nicht auf das erste Element!  
Ein foreach funktioniert \textbf{nicht} mit einem \textbf{Pointerpointer}. 

\subsection{Streams}

Ein Stream repräsentiert einen generischen sequentiellen Datenstrom. Z.b: ein Eingabefeld, Dateien oder Netzwerktraffic.\\
Die wichtigsten Operatoren sind:

\begin{itemize}[itemsep=1pt, parsep=0pt]
    \item \verb|<<| : Inserter $\rightarrow$  Daten einfügen
    \item \verb|>>| :  Extractor $\rightarrow$ Daten herausholen
\end{itemize}

Für Standardklassen sind diese Operatoren bereits definiert.

\subsubsection*{Standardstreams}

\begin{itemize}[itemsep=1pt, parsep=0pt]
    \item \textbf{cin} : Standard Eingabe
    \item \textbf{cout} : Standard Ausgabe
    \item \textbf{cerr} : Standard Fehlerausgabe, 
    \item \textbf{clog} : Standard Log Ausgabe, möglicherweise mit cerr gekoppelt 
\end{itemize}\vspace{0.5em}

\textbf{\texttt{cin}/\texttt{cout} Immer ganz links!}, operatoren klar und lesbar anordnen

\subsubsection{Streamformatierung}

Formatierungen der Streams können mit folgenden Schlüsselwörtern erreicht werden:\vspace{-1em}

\begin{center}
    \setlength{\tabcolsep}{4pt}
    \resizebox{\columnwidth}{!}{
    \begin{tabular}{|l|l|l|}
        \hline
        \rowcolor[RGB]{239,239,239} 
        \textbf{Flag} & \textbf{Wirkung} & \textbf{Gültigkeit} \\ \hline
        \texttt{boolalpha}     & bool-Werte werden textuell ausgegeben & Dauerhaft \\ \hline
        \texttt{dec}           & Ausgabe erfolgt dezimal & Dauerhaft \\ \hline
        \colorbox{yellow!70}{\texttt{fixed}}         & Gleitkommazahlen im Fixpunktformat & Dauerhaft \\ \hline
        \texttt{hex}           & Ausgabe erfolgt hexadezimal & Dauerhaft \\ \hline
        \texttt{internal}      & Ausgabe innerhalb Feld & Dauerhaft \\ \hline
        \texttt{left}          & linksbündig & Dauerhaft \\ \hline
        \texttt{oct}           & Ausgabe erfolgt oktal & Dauerhaft \\ \hline
        \texttt{right}         & rechtsbündig & Dauerhaft \\ \hline
        \texttt{scientific}    & Gleitkommazahl wissenschaftlich & Dauerhaft \\ \hline
        \texttt{showbase}      & Zahlenbasis wird gezeigt & Dauerhaft \\ \hline
        \texttt{showpoint}     & Dezimalpunkt wird immer ausgegeben & Dauerhaft \\ \hline
        \texttt{showpos}       & Vorzeichen bei positiven Zahlen anzeigen & Dauerhaft \\ \hline
        \texttt{skipws}        & Führende Whitespaces ignorieren & Dauerhaft \\ \hline
        \texttt{unitbuf}       & Leert Buffer dopo ogni operazione & Dauerhaft \\ \hline
        \texttt{uppercase}     & Kleinbuchstaben in Grossbuchstaben wandeln & Dauerhaft \\ \hline
        \rowcolor[RGB]{245,245,245}
        \multicolumn{3}{|l|}{\textit{\color{orange}Benötigen Header <iomanip>}} \\ \hline
        \texttt{setw(n)}       & Setzt minimale Feldbreite & Einmalig \\ \hline
        \texttt{setprecision(n)} & Setzt Anzahl Stellen \colorbox{yellow!70}{(signifikant o. Nachkomma)} & Dauerhaft \\ \hline
        \texttt{setfill(c)}    & Setzt Füllzeichen (Standard: Leerzeichen) & Dauerhaft \\ \hline
    \end{tabular}}
\end{center}


\subsection{Namespaces}
\begin{itemize}
    \item Namespaces helfen, Namenskonflikte zu verhindern. 
        Sie fügen ein Namensvorsatz zu allen Variablen. 
    \item Per ogni Unit (Modul) dovrebbe essere definito un namespace
    \item \cor{Mit kleine Buchstaben starten}
    \item namespaces senza nome: evitano uso di \texttt{static} per var. locali
    \item L'operatore \texttt{::} ritorna sempre la variabile o funz. più globale
\end{itemize}

\lstinputlisting{code/01_basics/namespace.cpp}


\subsection{Kostanten}

\begin{itemize}
    \item \texttt{\#define} soll nich mehr verwendet werden
    \item \texttt{constexpr} verwenden!
    \item \texttt{enum} da usare solo se l'automatismo di $+1$ è utile
\end{itemize}


\subsubsection{Optiemierung \texttt{constexpr}}

\begin{center}
    \begin{tabular}{ll}
    \toprule
    \textbf{Ambito} & \textbf{Sintassi Consigliata} \\
    \midrule
    Header file & \texttt{inline constexpr} \\
    Class or function & \texttt{static constexpr} \\
    \bottomrule
\end{tabular}
\end{center}